% SPDX-License-Identifier: Apache-2.0
\section{Recommendation}
\label{sec:recommendation}

\subsection{WebAuthn: a CTAPHID relay with a usbfs key side}

We recommend approach B. The remote half is fido-forward's sink, unchanged.
The Chromebook half is a new, small program, \code{ffwd-usbfs}, that reads
the key through usbfs and speaks fido-forward's wire format
(Section~\ref{sec:component}). Figure~\ref{fig:flow} shows the data flow.

\begin{figure*}[t]
\centering
\begin{tikzpicture}[node distance=4mm and 7mm]
  \node[mem] (yk) {YubiKey};
  \node[sbox,right=of yk,text width=24mm] (pb) {ChromeOS\\\code{permission\_broker}};
  \node[act,right=of pb,text width=30mm] (src) {\code{ffwd-usbfs}\\claims FIDO interface,\\refuses unsafe commands};
  \node[sbox,right=of src] (ssh) {\code{ssh -T}};
  \node[sbox,right=20mm of ssh] (sshd) {\code{sshd}};
  \node[act,right=of sshd,text width=26mm] (sink) {\code{fido-forward}\\\code{--sink}, setuid\\\code{fidofwd}};
  \node[mem,below=7mm of sink,text width=26mm] (hid) {virtual \code{/dev/hidrawM}\\page 0xF1D0};
  \node[sbox,left=of hid,text width=22mm] (chr) {Chrome in the CRD session};
  \draw[bidi] (yk) -- node[above,font=\scriptsize]{USB} (pb);
  \draw[bidi] (pb) -- node[above,font=\scriptsize]{usbfs} (src);
  \draw[bidi] (src) -- node[above,font=\scriptsize]{stdio} (ssh);
  \draw[bidi,double] (ssh) -- node[above,font=\scriptsize]{SSH} (sshd);
  \draw[bidi] (sshd) -- node[above,font=\scriptsize]{stdio} (sink);
  \draw[bidi] (sink) -- node[right,font=\scriptsize]{\code{/dev/uhid}} (hid);
  \draw[bidi] (hid) -- (chr);
  \begin{scope}[on background layer]
    \node[draw,dashed,rounded corners,fit=(yk)(pb)(src)(ssh),inner sep=2mm,
          label={[font=\scriptsize]above:Chromebook (Crostini: \code{ffwd-usbfs}, \code{ssh})}] {};
    \node[draw,dashed,rounded corners,fit=(sshd)(sink)(hid)(chr),inner sep=2mm,
          label={[font=\scriptsize]above:Remote host}] {};
  \end{scope}
\end{tikzpicture}
\caption{Data flow of the recommended design. Only the FIDO interface leaves
the Chromebook. SSH is the only transport and the only authentication; no
TCP port opens on either host.}
\label{fig:flow}
\end{figure*}

\begin{example}[One login]
\label{ex:login}
\begin{enumerate}
\item Chrome on the remote host writes a \code{CTAPHID\_CBOR} request with
      \code{getAssertion} to \code{/dev/hidrawM}.
\item The sink reads it from \code{/dev/uhid} and sends the 64-byte reports
      up the SSH channel.
\item \code{ffwd-usbfs} joins the reports, checks the command against its
      policy, and writes them to the key's interrupt OUT endpoint.
\item The key sends keepalives until the user touches it, then the
      response.
\item \code{ffwd-usbfs} reads them from the interrupt IN endpoint and sends
      them down the SSH channel.
\item The sink writes them to \code{/dev/uhid}. Chrome reads the response
      and completes the login.
\end{enumerate}
\end{example}

Five reasons support this choice:
\begin{enumerate}
\item Only the FIDO interface leaves the Chromebook. PIV, OpenPGP and OTP
      stay local.
\item The relay refuses commands that erase or change the key
      (Section~\ref{sec:security}).
\item The remote half exists, drops privilege before it reads input, and
      builds with Bazel here for amd64 and arm64.
\item SSH is the only transport and authentication.
\item Keepalives pass through, so a slow touch does not time out.
\end{enumerate}

\subsection{OpenPGP and PIV: agent forwarding}

For OpenPGP, forward gpg-agent's extra socket with \code{RemoteForward}, and
run \code{scdaemon} and \code{pcscd} in Crostini~\cite{gnupgfwd}. For PIV,
run \code{p11-kit server} with the YubiKey PKCS\#11 module in Crostini and
forward its socket~\cite{p11kit}. The FIDO relay claims only the FIDO
interface, and \code{pcscd} claims only the CCID interface, so the two should
work side by side (assumption; test T6).

FIDO keys for SSH cannot be used from Crostini (Finding~\ref{find:nohidraw}).
On the remote host they work through the virtual key of approach B, because
the remote host has a \code{hidraw} node.

\subsection{Fallback: USB/IP with usbipdcpp}

Use approach A if the probe fails, or if you need PIV and OpenPGP as real
devices on the remote host. Use usbipdcpp, not jiegec/usbip. Use it only with
a remote host that you trust fully, because it gives that host's root user
full control of the key.
